\documentclass[10pt]{article}
\usepackage[margin=0.60in]{geometry}
\usepackage{booktabs,enumitem,fancyhdr,graphicx,listings,tabularx,titlesec,xcolor}
\usepackage[hidelinks]{hyperref}
\graphicspath{{docs/lessons/assets/}{assets/}}
\hypersetup{
 pdftitle={ADK Lesson 064: Owned Single-Wire Transactions},
 pdfauthor={Kevin Shanahan},
 pdfsubject={A bounded copied intent and receipt policy for typed DS18B20-family single-wire transactions},
 pdflang={en-US}}
\pagestyle{fancy}\fancyhf{}
\lhead{\textsc{ADK Lesson 064}}\rhead{Owned Single-Wire Transactions}
\cfoot{\thepage}\setlength{\headheight}{14pt}
\setlength{\parindent}{0pt}\setlength{\parskip}{3pt}
\setlist{nosep,leftmargin=1.45em}
\titleformat{\section}{\Large\bfseries}{\thesection}{0.5em}{}
\newcommand{\lessonbox}[1]{\begin{center}\fbox{\parbox{0.92\linewidth}{#1}}\end{center}}
\newcommand{\blankline}{\rule{0.76in}{0.4pt}}
\newcommand{\longblank}{\rule{1.35in}{0.4pt}}
\lstset{
 basicstyle=\ttfamily\footnotesize,
 columns=fullflexible,
 breaklines=true,
 frame=single,
 rulecolor=\color{black},
 showstringspaces=false}

\begin{document}
\small
\begin{center}
 {\Huge\bfseries Owned Single-Wire Transactions}\\[3pt]
 {\Large Prove a bounded policy before touching a bus}\\[7pt]
 \fbox{\textbf{E0 SYNTHETIC REPLAY --- NO PIN, PULL-UP, OR SENSOR}}
\end{center}
\begin{tabularx}{\linewidth}{@{}>{\bfseries}l X >{\bfseries}l X@{}}
Lesson & 064 --- component & Time & 180--240 minutes\\
Status & host verified and published at E0 & Board & compile-only Mega replay\\
Input & copied requests and line receipts & Observation & named intent/snapshot cells\\
\end{tabularx}

\lessonbox{\textbf{Responsibility.}
\texttt{OneWireTransactionPolicy} expands one reviewed typed operation into a
bounded sequence of copied line intents, validates one correlated copied
receipt at a time, retains an attributable accepted prefix, and requests
release during completion or rollback. It cannot drive, sample, power, or
release a physical line.}

\section{Separate protocol policy from electrical execution}
The DS18B20 data sheet describes an open-drain 1-Wire interface, reset and
presence behavior, read and write slots, ROM commands, supply modes, and
device timing limits.\footnote{\href{https://www.analog.com/media/en/technical-documentation/data-sheets/ds18b20.pdf}{Analog Devices, \emph{DS18B20 Programmable Resolution 1-Wire Digital Thermometer}, Rev. 6.}}
Those facts explain what a future qualified adapter must implement. They do
not prove that an unidentified kit part, a Mega pin, or a particular pull-up
meets the requirements.

Lesson 064 therefore stops at a copied evidence boundary. The policy emits
semantic values such as \texttt{DriveLow}, \texttt{Release}, and
\texttt{Sample}. A synthetic producer returns a copied receipt saying what it
applied or observed. E0 proves only that the policy processes that evidence
deterministically.

\begin{center}
% ADK visual: pencil
\includegraphics[width=0.96\linewidth,height=1.85in,keepaspectratio]
 {064-evidence-boundary-pencil.png}

\small\textit{Alternative text: a graphite policy box hands two blank copied
cards toward a locked gate; a faint generic sensor silhouette remains isolated
beyond the gate, and no physical connection crosses it.}
\end{center}

\begin{tabularx}{\linewidth}{@{}X X@{}}
\toprule
E0 can establish & E0 cannot establish\\
\midrule
typed phase order and receipt correlation & pin ownership, voltage, or open-drain behavior\\
configured timing-window acceptance & physical pulse width, rise time, or interrupt margin\\
bounded rollback and release intent & a physical line reached idle\\
byte-stable replay from copied evidence & genuine DS18B20 identity or thermal behavior\\
\bottomrule
\end{tabularx}

There is no formal electrical schematic or wiring exercise in this lesson.
An authoritative schematic requires an exact specimen, package and pinout,
ordinary pull-up and supply choice, endpoint ownership, measured rails and
waveforms, and separate resource and safe-state evidence.

\section{Use a closed operation vocabulary}
The public request chooses one reviewed operation. Reset/presence, one bounded
ROM-search pass, and single-drop ROM read cover identification mechanics.
Addressed operations cover power-supply read, conversion start, one conversion
status continuation, and scratchpad read. Their public enum spellings are
\texttt{ResetPresence}, \texttt{SearchRomPass},
\texttt{ReadRomSingleDrop}, the four \texttt{MatchRom...} operations. No
request accepts arbitrary command bytes, captured pulse trains, or an unknown
replay script.

\begin{center}
% ADK visual: pencil
\includegraphics[width=0.96\linewidth,height=1.75in,keepaspectratio]
 {064-typed-operations-pencil.png}

\small\textit{Alternative text: four graphite index cards progress through a
reset gesture, an identity tag, bounded slot marks, and an open-hand release
gesture; the surrounding prose supplies the operation names.}
\end{center}

Addressed operations reject only the all-zero ROM. Every other complete
64-bit value is transported byte-for-byte, including a wrong-family or
bad-CRC value. Lesson 064 does not call it valid. Lesson 065 exclusively owns
family-\texttt{0x28} and Dallas CRC qualification.

\texttt{ReadRomSingleDrop} is admitted only when configuration explicitly
names a single-drop fixture. Multidrop work uses exactly one bounded
\texttt{SearchRomPass}. The manufacturer search algorithm explains why each
device's 64-bit ROM participates in a discrepancy tree,\footnote{\href{https://www.analog.com/en/resources/app-notes/1wire-search-algorithm.html}{Analog Devices, \emph{1-Wire Search Algorithm}, Application Note 187.}}
but its enumeration order is evidence order, never physical position.

Every ordinary operation begins with reset and presence.
\texttt{ResetPresence} stops there without sending a command. The sole reset
exception is \texttt{MatchRomReadConversionStatus}: it is admitted only as a
one-slot continuation of this policy's immediately preceding completed
\texttt{MatchRomStartConversion} for the same nonzero ROM. A free-standing
status request rejects atomically.

Why should a wrong-family nonzero ROM be transported rather than ``fixed'' by
Lesson 064? \longblank.

\section{Make microsecond time explicit}
The repository's millisecond time type cannot represent single-wire slot
windows. Lesson 064 uses a dedicated unsigned 32-bit
\texttt{MicrosecondTimePoint} and \texttt{MicrosecondDuration}. Configuration
names independent lower and upper limits for reset low, reset release,
presence start, presence low, write-zero low, write-one low, read initiation,
read sampling, complete slot, recovery, and the whole transaction deadline.

\begin{center}
% ADK visual: pencil
\includegraphics[width=0.96\linewidth,height=1.70in,keepaspectratio]
 {064-microsecond-window-pencil.png}

\small\textit{Alternative text: a graphite timeline on torn paper has four
ticks and a shaded admissible window between the middle ticks; a separate
clock card represents caller-supplied microsecond time.}
\end{center}

Tests inspect immediately below, exactly at, and immediately above every lower
and upper boundary. All configured durations stay strictly below the modular
half range. Backward time and an exact half-range difference are ambiguous
and reject atomically; an ordinary unsigned rollover with a valid forward
distance is accepted.

Analog Devices' software guidance emphasizes that recovery and sample timing
depend on the device and network conditions.\footnote{\href{https://www.analog.com/en/resources/design-notes/1wirereg-extended-network-standard.html}{Analog Devices, \emph{1-Wire Extended Network Standard}, discussing Application Note 126 and network-dependent timing.}}
The configured E0 windows are copied-evidence acceptance rules, not a claim
that any physical waveform met them.

\section{Correlate every intent and receipt}
An operation request carries request sequence, typed operation, optional
addressed ROM or fixed search state, supplied start time, supply mode, and
producer status. Every emitted intent binds owner, lifecycle and configuration
revisions, request and transaction generations, phase sequence, slot index,
operation, line intent, legal action window, and addressed ROM.

A receipt carries producer source and configuration, its own sequence and
observation time, then repeats owner, lifecycle, request, transaction,
operation, phase, phase sequence, slot, and applied-line correlation. It adds
the copied sampled level, acceptance bit, and producer status. The addressed
ROM and legal window remain in the owed intent and are checked through that
binding; they are not receipt fields.

\begin{center}
% ADK visual: pencil
\includegraphics[width=0.96\linewidth,height=1.65in,keepaspectratio]
 {064-intent-receipt-pencil.png}

\small\textit{Alternative text: two graphite cards face each other with the
same ordered set of geometric seals and opposing arrows, representing exact
field correlation and one return receipt.}
\end{center}

A matching receipt advances at most one phase or slot. A byte-identical
duplicate is idempotent. Changed duplicates, crossed transactions, a foreign
source, wrong generation, wrong operation, or wrong phase or slot reject
atomically. A structurally matching receipt outside its inclusive window is
retained as attributed timeout evidence and starts one release rollback; it
does not advance the requested operation.

Sequence regression and half-range ambiguity reject atomically. If a call
cannot leave a unique next transaction, phase, lifecycle, receipt, or cleanup
correlation, the entire call returns \texttt{CapacityExceeded} without
mutation. No counter wraps to zero.

\begin{enumerate}
 \item If every field matches except the slot index, should the policy use the
       sampled bit? \blankline.
 \item If a byte-identical receipt is presented twice, how many phase
       transitions occur in total? \blankline.
 \item A receipt has valid correlation but producer failure. Does structure
       reject, or does rollback begin?\\[2pt] \longblank.
\end{enumerate}

\section{Serialize least-significant bits first}
The DS18B20 protocol sends data and commands least-significant bit first.
Lesson 064 freezes the request at \texttt{begin()}, expands the allowlisted
operation into bounded private slot mechanics, and assembles copied read bits
back into bytes in the same order.

\begin{center}
% ADK visual: pencil
\includegraphics[width=0.96\linewidth,height=1.65in,keepaspectratio]
 {064-lsb-prefix-pencil.png}

\small\textit{Alternative text: eight graphite cards with progressively
darkened dots run left to right into a folder, representing least-significant
first slot order and an attributable accumulated prefix.}
\end{center}

Partial read bytes are prefix evidence, not a completed device response. A
failure at the first, middle, or final slot preserves exactly the admitted
prefix and moves to cleanup. Neither an incomplete scratchpad nor a copied
search branch becomes a temperature.

For copied bits \(1,0,1,1,0,0,0,1\) in slot order, write the assembled byte
in binary: \blankline. Which bit becomes the most significant bit?
\blankline.

\section{Treat release as an owed transaction}
The policy never emits ``drive high.'' Its canonical inactive and cleanup
line intent is \texttt{Release}. Initialization, reset, cancellation, timeout,
producer failure, and shutdown may all leave one correlated release receipt
owed. Exhaustion is the exception: if no unique cleanup correlation can be
minted, the entire triggering call rejects atomically rather than promising
an unconfirmable release.

\begin{center}
% ADK visual: pencil

\includegraphics[width=0.96\linewidth,height=2.05in,keepaspectratio]
 {064-release-cleanup-pencil.png}

\small\textit{Alternative text: four separated graphite trays show active
work, placement of an open-hand release card, a pending stack, and a confirmed
card; a failed card lies separately below and does not overlap the path.}
\end{center}

\begin{tabularx}{\linewidth}{@{}l X@{}}
\toprule
Evidence & Permitted statement\\
\midrule
policy emitted \texttt{Release} & software requested line release\\
matching copied release receipt & producer reported applying the request\\
future E1 waveform observation & qualified physical line reached its idle window\\
\bottomrule
\end{tabularx}

\texttt{confirmCleanup()} is the sole closing seam. Reset becomes ready only
after its matching copied release receipt. Shutdown remains initialized in a
closing state so it can accept that receipt, then becomes inert. Shutdown
supersedes any earlier pending cleanup with a fresh lifecycle and correlated
\texttt{Release}; the earlier receipt becomes stale. Without the closing
receipt, the snapshot remains \texttt{ReleaseUnconfirmed}; E0 never says the
physical bus is idle.

\texttt{completedEvidence()} returns \texttt{ResourceBusy} after successful
data slots while terminal release remains unconfirmed. It returns
\texttt{Ok} with the complete attributed snapshot only after that receipt,
and returns \texttt{NotInitialized} for every other state.

\section{Advance deadlines without polling}
Ordinary phase transitions require one receipt through \texttt{update()}.
The receipt-free \texttt{advance(now,intent)} path does only one job: inspect
supplied time. At or before the deadline it leaves the owed phase unchanged.
Only the first supplied time strictly after the deadline emits at most one
correlated release intent. It never polls, retries, catches up, invents phase
evidence, or synthesizes a cleanup receipt.

\begin{center}
% ADK visual: pencil
\includegraphics[width=0.96\linewidth,height=1.85in,keepaspectratio]
 {064-deadline-cleanup-pencil.png}

\small\textit{Alternative text: a graphite hourglass runs, then empties beside
one open-hand release card, which remains in a tray as the same unresolved
cleanup obligation on later calls.}
\end{center}

\texttt{MatchRomReadConversionStatus} is a one-slot continuation, not a
free-standing poll transaction. It is available only after this policy
completed \texttt{MatchRomStartConversion} for the same ROM, and it reports
one copied bit. A later status check requires another completed conversion
start; the component hides no poll loop.

Predict the result of calling \texttt{advance()} repeatedly after a timeout
but before cleanup confirmation: number of successful phase receipts invented
\blankline; cleanup quality \longblank.

\section{Keep physical external-power evidence open}
Only the copied \texttt{ExternallyPowered} supply-mode value is admitted at
E0; admitting that enum does not energize or qualify a circuit.
\texttt{ParasitePower} returns \texttt{Unsupported} before state changes.
There is no strong-pull-up phase or intent. Ordinary \texttt{Release} must
never be reinterpreted as output-high drive or strong-pull power.

The data sheet explains that parasite operation uses the data line for power
and imposes additional strong-pull requirements during operations such as
temperature conversion. That is precisely why it is excluded here. A future
parasite-power endpoint needs separate topology, current, rail, timing,
exclusivity, rollback, stuck-on handling, and independent power-removal
evidence.

\begin{tabularx}{\linewidth}{@{}l l X@{}}
\toprule
Request & E0 disposition & Reason\\
\midrule
externally powered typed operation & eligible & copied protocol policy only\\
parasite-powered operation & \texttt{Unsupported} & no reviewed power capability\\
arbitrary captured transaction & invalid & outside closed operation surface\\
scratchpad EEPROM copy & absent & endurance and strong-pull work excluded\\
\bottomrule
\end{tabularx}

\section{Replay the state machine}
The compile-only Mega sketch supplies microsecond values, requests, receipts,
and fixed caller-owned output buffers. It does not include a GPIO adapter and
must not be uploaded as a bus driver.

\begin{lstlisting}
OneWireStepIntent intent = {};
OneWireTransactionSnapshot snapshot = {};
OneWireTransactionSnapshot evidence = {};

OneWireTransactionPolicy transaction(config);
Status status = transaction.initialize(now, intent);
status = transaction.confirmCleanup(now, releaseReceipt);
status = transaction.begin(now, request, intent);
status = transaction.update(now, phaseReceipt, intent);
status = transaction.snapshot(snapshot);
status = transaction.completedEvidence(evidence);
\end{lstlisting}

\begin{center}
% ADK visual: pencil
\includegraphics[width=0.96\linewidth,height=1.80in,keepaspectratio]
 {064-replay-ledger-pencil.png}

\small\textit{Alternative text: a stack of graphite request cards and a
separate receipt stack enter an abstract policy box, which fills five blank
result cells in a row for the replay ledger.}
\end{center}

The canonical sketch records only baseline cleanup confirmation, parasite
rejection, one reset/presence transaction and its terminal release,
receipt-provenance predictions, one strictly-post-deadline timeout and
rollback confirmation, and shutdown confirmation. Its named result cells do
not claim to replay addressed reads, Search ROM, or malformed receipts.

The host proof matrix separately covers every timing neighbor; presence absent, valid, early,
late, and stuck; both copied read values; least-significant-first assembly;
first, middle, and final slot failures; fieldwise changed duplicates; crossed
correlation; single-drop restrictions; bounded search branches and collision;
parasite rejection; timeout; cancellation; reset and shutdown cleanup;
rollover and exhaustion; and two independent policies producing byte-stable
results from the same trace.

\section{Inspect gates before making a claim}
\begin{tabularx}{\linewidth}{@{}l X X@{}}
\toprule
Gate & E0 evidence & Still open\\
\midrule
resources & zero hardware owners and exact compiler report & physical pin/timer/interrupt ownership\\
protocol & typed phases and copied timing boundaries & measured pulse, sample, rise, and recovery timing\\
cleanup & correlated Release intent and receipt & physical high-impedance/idle voltage\\
identity & byte-exact nonzero ROM transport & family, CRC, specimen authenticity\\
power & parasite requests reject atomically & any strong-pull topology or campaign\\
\bottomrule
\end{tabularx}

The ordinary Mega replay measures 10,390 bytes of flash and 735 bytes of
static SRAM. The exact isolated no-LTO probe measures 12,116 bytes of flash,
735 bytes of static SRAM, and 199 bytes of conservative synchronous stack.
The policy object is 254 bytes: an independently accepted miss above the
192-byte target, subject to a zero-growth constraint, and only two bytes below
the 256-byte hard limit. Caller-owned request, search, intent, receipt, and
snapshot values measure 29, 10, 45, 39, and 84 bytes, or 207 bytes together.
After the specified reserve and stack accounting, 7,130 bytes of Mega SRAM
remain.

The ordinary sketch and exact isolated probe are separate measurement
boundaries; their flash and static-SRAM values must not be added. These
compiler measurements do not establish electrical timing, physical cleanup,
or sensor behavior. Independent review accepts both flash target misses and
the object target miss; every hard and residual gate passes. The stable
evidence fingerprint is
\texttt{77c3da45f8ea0f6943a01f57fdd2062a78a1fb61fc79837c1c15c45fbd876b97}.

\lessonbox{\textbf{E1a remains open.}
Before any powered adapter, qualify an exact externally powered specimen and
ordinary external pull-up. Record package and pinout, genuine/compatible
evidence, supply and current, drive-low and release rails, rise time,
reset/presence/slot waveforms, interrupt effects, timeout cleanup, stuck-low
behavior, startup/reset/shutdown, and an independent physical line-release
observation. Parasite power remains E1b-closed.}

\section{Completion record}
\begin{tabularx}{\linewidth}{@{}X X@{}}
\toprule
Check & Result or evidence\\
\midrule
Closed typed operation and nonzero-ROM transport boundary & \\
\addlinespace[1.1em]
Every configured microsecond-window neighbor and rollover & \\
\addlinespace[1.1em]
Intent/receipt correlation, duplicate, regression, and exhaustion & \\
\addlinespace[1.1em]
Least-significant-first slots and attributable partial prefix & \\
\addlinespace[1.1em]
Timeout, cancellation, reset, shutdown, and cleanup confirmation & \\
\addlinespace[1.1em]
Parasite rejection and absence of hidden polling/retry & \\
\addlinespace[1.1em]
Canonical replay, exact resource evidence, and open E1 gates & \\
\addlinespace[1.1em]
\bottomrule
\end{tabularx}

\begin{center}
\fbox{\parbox{0.90\linewidth}{
\textbf{Completion statement:} I can distinguish copied protocol evidence
from electrical execution, correlate every bounded phase, preserve partial
prefix and cleanup debt, explain why release is not proof of an idle bus, and
name the physical evidence still required.

\vspace{0.10in}
Learner: \longblank\quad Reviewer: \longblank\quad Date: \blankline
}}
\end{center}

\section*{Primary references}
\begin{itemize}
 \item Analog Devices, \emph{DS18B20 Programmable Resolution 1-Wire Digital
       Thermometer}, Rev. 6,
       \href{https://www.analog.com/media/en/technical-documentation/data-sheets/ds18b20.pdf}{manufacturer data-sheet PDF}.
 \item Analog Devices, \emph{1-Wire Search Algorithm}, Application Note 187:
       \href{https://www.analog.com/en/resources/app-notes/1wire-search-algorithm.html}{manufacturer application-note page}.
 \item Analog Devices, \emph{1-Wire Extended Network Standard}, including its
       discussion of Application Note 126:
       \href{https://www.analog.com/en/resources/design-notes/1wirereg-extended-network-standard.html}{manufacturer design-note page}.
\end{itemize}

This print edition is untagged. Its text is searchable and copyable, and the
HTML lesson provides the hardware-independent alternative. No PDF/UA, WCAG,
or other PDF accessibility-conformance claim is made.

\end{document}
